\section{Discussion and Conclusions}

\subsection{Discussion of the Results}

The work shows that an important part of the cost of declarative dynamic heuristics can be moved from grounding to search. On BSP, native \mintinline{text}|#heuristic| directives grow as $\Theta(n^3)$ and exceed one million ground objects at $n=100$; the lazy variants keep two facts, and the difference is visible in the propositional instance size, in the peak memory, and in the timeout profile. The mechanism generalizes: whenever the heuristic weight or applicability depends on an aggregate over choice atoms, the grounder must enumerate the value domain of the aggregate, and the same product structure that explodes BSP also drives the PUP heuristics, whose ground form multiplies four unrolled value domains per ground pair.

This result should not be read as a replacement for general lazy grounding. clingo still grounds the base program; the propagator acts only on the heuristic component. Nor should it be read as a claim that laziness is free: the runtime work is real, it is paid at every decision, and the per-decision metrics exist precisely to keep that cost on the record. The honest formulation is a space--time trade-off with a crossover: below it, ground-and-solve remains competitive, because the explosion is manageable and the lazy machinery (in the Prolog case, an entire in-process logic-programming engine) carries a fixed overhead; beyond it, the constant-size heuristic component is the only representation that survives.

The second central finding is semantic, and it is a finding about \emph{meaning}, not performance. The same declarative heuristic changes behaviour dramatically depending on how partial information is read. Under the Alpha-like reading, the BSP heuristic is a perfect greedy strategy: $n$ choices, zero conflicts. Under the clingo-like reading, the very same weights and bodies barely move the search away from the baseline, because the negative guards open only on established falsity and the aggregates wait for their sources to be determined. Neither reading is \enquote{wrong}: they answer different questions, and the four-variant matrix exists to keep those questions separate. A comparison that silently mixed the axes, for example by measuring \mintinline{text}|la| against \mintinline{text}|gc| and attributing the difference to laziness, would be methodologically invalid; the thesis takes some care never to do so.

The third finding concerns the translation layer between the two heuristic languages. Alpha's weight--priority pairs are global levels; clingo's priorities arbitrate on a single atom. Every encoding that expresses an ordering between different atom families must therefore flatten its levels into weights wherever clingo's reading applies: in the ground variants, because clasp is the consumer, and in the lazy clingo variants, because the propagator deliberately reproduces clingo's behaviour there. This is easy to state and easy to get wrong: an early revision of the lazy-clingo PUP and HRP encodings carried Alpha-style priorities that the clingo semantics does not rank globally, producing an unintended decision order that inverted the designed one (sensors before zones; cabinet-filling before reuse), and diverging from the Prolog twin encodings of the same variants. The final audit caught and corrected the discrepancy by applying the same flattening used in the ground encodings. The episode is instructive: the interaction between priority semantics and evaluation backend is exactly the kind of silent, non-crashing behavioural divergence that a benchmark pipeline cannot detect from exit codes, and it motivates the equivalence and wiring guards that the suite now includes.

A related engineering lesson comes from the silent-fallback failure mode of the Prolog backend. A single missing environment variable degrades every lazy Prolog run to the no-heuristic baseline without any error, since the fallback evaluator simply fails on the Prolog-style rule bodies. The pipeline defends itself with positive evidence of activity (\mintinline{text}|decide_calls > 0|), cross-backend agreement checks, and a wrapper that forces the variable; the general principle, that a heuristic mechanism must prove it ran rather than merely not fail, seems worth stating explicitly.

Finally, the auxiliary-form experiment delimits the approach from within. \mintinline{text}|la_aux| shows that laziness only pays if the \emph{entire} dynamic part of the heuristic, aggregate included, stays out of the grounder: materializing the heuristic body into an auxiliary relation reintroduces the very product the mechanism was built to avoid.

\subsection{Limitations}

The current implementation has explicit technical limitations. The native backend indexes targets and bodies through numeric tuples only; non-numeric arguments are not supported. The arithmetic language covers addition, subtraction, and multiplication; the available dynamic aggregates are sum, count, min, and max, over a single source predicate, with joins delegated to precompiled auxiliary predicates in the encoding. The modifiers \mintinline{text}|init| and \mintinline{text}|factor| are recognized but deliberately rejected: in native clingo they modify solver-internal scores that the \mintinline{text}|decide| interface cannot reach, and approximating them as \mintinline{text}|level| would silently change their meaning. The propagator is registered sequentially and has no synchronization layer for parallel solving. The tie-break between targets of equal rank is deterministic but does not reproduce clingo's internal scores.

On the evaluation side, the completed quantitative analysis concerns BSP; the PUP and HRP campaigns are wired and validated but their definitive figures await the re-run with the corrected lazy-clingo encodings. The memory metric is end-to-end process RSS, which is the honest but blunt instrument discussed in Sections~2 and~3: it includes the Prolog engine for the lazy Prolog variants and cannot separate grounding memory from search memory. The auxiliary clingo evaluator of the Prolog build re-grounds a program per decision and is a functional reference, not a performance-comparable backend. Finally, the whitelist of static predicates and the inference of the program constant $n$ in the Prolog backend are benchmark-specific conveniences that a production version should replace with a declarative interface.

\subsection{Implications}

The main implication is that declarative heuristics can be treated as an intermediate layer between modelling and solving. One does not have to choose between fully grounded \mintinline{text}|#heuristic| directives and procedural heuristics hard-coded in the solver: a compact declarative surface, evaluated at runtime with access to the partial assignment, is a workable third point in the design space, and it can host more than one evaluation semantics behind one syntax. The two backends demonstrate two implementation strategies for that layer, a compiled, incremental one and an embedded general-purpose engine, with the expected trade-off between per-decision efficiency and expressive freedom of the rule bodies.

From an engineering perspective, the work also shows that clingo's propagator and heuristic interfaces are flexible enough to prototype non-trivial guidance mechanisms without touching the solver core: watched literals, incremental aggregate states, per-target modifier resolution, and a global decision ranking are all expressible against the public API of \mintinline{text}|Clingo::Heuristic|.

\subsection{Future Work}

Several extensions follow naturally. The mechanism should become configurable rather than always-on, and the expression language deserves integer division, modulo, and comparisons, together with a principled treatment of non-numeric symbols through a stable internal mapping. Aggregates with in-template joins would remove the need for precompiled auxiliary predicates. On the Prolog side, the synchronization protocol could move from per-literal goals to batched updates, and the backend could expose its per-decision statistics through clingo's statistics interface instead of stderr parsing. A parallel-solving synchronization layer and a tie-break policy that cooperates with the solver's own scores are natural solver-side improvements.

On the evaluation side, the immediate step is the completed PUP and HRP campaigns with the corrected encodings, including the cross-backend comparison of the per-decision costs on aggregate-heavy (PUP) versus negation-heavy (HRP) heuristics. Beyond that, configuration, scheduling, and packing domains are natural candidates, precisely because their useful heuristics are dynamic. A further direction is a verified converter from annotated \mintinline{text}|#heuristic| directives to lazy facts, with the flattening rule of Section~2 applied mechanically rather than by hand; the audit episode discussed above is a concrete argument for automating that step.

\subsection{Conclusions}

This thesis presented a lazy evaluation mechanism for declarative domain-specific heuristics in clingo, instantiated in two backends over a common propagator design: a native C++ backend interpreting compact \mintinline{text}|__heuristic| facts with incremental dynamic aggregates, and an in-process SWI-Prolog backend evaluating heuristic rules as logic programs against a synchronized image of the solver trail. The mechanism supports two explicit readings of partial information, the clingo-like and the Alpha-like semantics, and reproduces clingo's weight, priority, and modifier behaviour where that reading applies, including the systematic flattening of Alpha's global priority levels into clingo weights.

The evaluation on BSP shows a reduction of the ground heuristic component from $\Theta(n^3)$ directives, more than one million at $n=100$, to two constant facts, with the corresponding asymptotic saving in real memory, and quantifies the price: a per-decision runtime cost that ground-and-solve does not pay. The suite validates that all variants and both backends preserve the answer sets, guards against silent misconfiguration, and extends to PUP and HRP.

The final result is a working, documented, and measurable prototype together with a methodology: two orthogonal axes that keep representation and semantics separate, a flattening rule that makes Alpha-style encodings meaningful under clingo's reading of priorities, and an experimental discipline that states what each number measures. It does not solve the grounding bottleneck in general, but it provides a concrete, verifiable path for bringing dynamic declarative heuristics into clingo without always paying the cost of their full ground materialization.
